\chapter{2021年全国乙卷（文）}

\section{单选题}

\begin{problem}
已知全集 \(U=\{1,2,3,4,5\}\)，集合 \(M=\{1,2\}\)，\(N=\{3,4\}\)，则 \(C_U\left(M\cup N\right)=\)（\quad）

\choices
    {\(\{5\}\)}
    {\(\{1,2\}\)}
    {\(\{3,4\}\)}
    {\(\{1,2,3,4\}\)}
\end{problem}

\begin{answer}
A
\end{answer}

\begin{problem}
设 \(\i z=4+3\i\)，则 \(z=\)（\quad）

\choices
    {\(-3-4\i\)}
    {\(-3+4\i\)}
    {\(3-4\i\)}
    {\(3+4\i\)}
\end{problem}

\begin{answer}
C
\end{answer}

\begin{problem}
已知命题 \(p:\exists x\in\R\)，\(\sin x<1\)；命题 \(q:\forall x\in\R\)，\(\e^{|x|}\ge1\)，则下列命题中为真命题的是（\quad）

\choices
    {\(p\wedge q\)}
    {\(\neg p\wedge q\)}
    {\(p\wedge \neg q\)}
    {\(\neg\left(p\vee q\right)\)}
\end{problem}

\begin{answer}
A
\end{answer}

\begin{problem}
函数 \(f(x)=\sin\dfrac{x}{3}+\cos\dfrac{x}{3}\) 的最小正周期和最大值分别是（\quad）

\choices
    {\(3\pi\) 和 \(\sqrt2\)}
    {\(3\pi\) 和 \(2\)}
    {\(6\pi\) 和 \(\sqrt2\)}
    {\(6\pi\) 和 \(2\)}
\end{problem}

\begin{answer}
C
\end{answer}

\begin{problem}
若 \(x\)，\(y\) 满足约束条件
\(\begin{cases}
x+y\ge4,\\
x-y\le2,\\
y\le3,
\end{cases}\)，
则 \(z=3x+y\) 的最小值为（\quad）

\choices
    {\(18\)}
    {\(10\)}
    {\(6\)}
    {\(4\)}
\end{problem}

\begin{answer}
C
\end{answer}

\begin{problem}
\(\cos^2\dfrac{\pi}{12}-\cos^2\dfrac{5\pi}{12}=\)（\quad）

\choices
    {\(\dfrac{1}{2}\)}
    {\(\dfrac{\sqrt3}{3}\)}
    {\(\dfrac{\sqrt2}{2}\)}
    {\(\dfrac{\sqrt3}{2}\)}
\end{problem}

\begin{answer}
D
\end{answer}

\begin{problem}
在区间 \(\left(0,\dfrac{1}{2}\right)\) 随机取 \(1\) 个数，则取到的数小于 \(\dfrac{1}{3}\) 的概率为（\quad）

\choices
    {\(\dfrac{3}{4}\)}
    {\(\dfrac{2}{3}\)}
    {\(\dfrac{1}{3}\)}
    {\(\dfrac{1}{6}\)}
\end{problem}

\begin{answer}
B
\end{answer}

\begin{problem}
下列函数中最小值为 \(4\) 的是（\quad）

\choices
    {\(y=x^2+2x+4\)}
    {\(y=|\sin x|+\dfrac{4}{|\sin x|}\)}
    {\(y=2^x+2^{2-x}\)}
    {\(y=\ln x+\dfrac{4}{\ln x}\)}
\end{problem}

\begin{answer}
C
\end{answer}

\begin{problem}
设函数 \(f(x)=\dfrac{1-x}{1+x}\)，则下列函数中为奇函数的是（\quad）

\choices
    {\(f(x-1)-1\)}
    {\(f(x-1)+1\)}
    {\(f(x+1)-1\)}
    {\(f(x+1)+1\)}
\end{problem}

\begin{answer}
B
\end{answer}

\begin{problem}
在正方体 \(ABCD-A_1B_1C_1D_1\) 中，\(P\) 为 \(B_1D_1\) 的中点，则直线 \(PB\) 与 \(AD_1\) 所成的角为（\quad）

\choices
    {\(\dfrac{\pi}{2}\)}
    {\(\dfrac{\pi}{3}\)}
    {\(\dfrac{\pi}{4}\)}
    {\(\dfrac{\pi}{6}\)}
\end{problem}

\begin{answer}
D
\end{answer}

\begin{problem}
设 \(B\) 是椭圆 \(C:\dfrac{x^2}{5}+y^2=1\) 的上顶点，点 \(P\) 在 \(C\) 上，则 \(\lvert PB\rvert\) 的最大值为（\quad）

\choices
    {\(\dfrac{5}{2}\)}
    {\(\sqrt6\)}
    {\(\sqrt5\)}
    {2}
\end{problem}

\begin{answer}
A
\end{answer}

\begin{problem}
设 \(a\ne0\)，若 \(x=a\) 是函数 \(f(x)=a\left(x-a\right)^2\left(x-b\right)\) 的极大值点，则（\quad）

\choices
    {\(a<b\)}
    {\(a>b\)}
    {\(ab<a^2\)}
    {\(ab>a^2\)}
\end{problem}

\begin{answer}
D
\end{answer}

\section{填空题}

\begin{problem}
已知向量 \(\bs{a}=(2,5)\)，\(\bs{b}=(\lambda,4)\)，若 \(\bs{a}\parallel\bs{b}\)，则 \(\lambda=\)\fillinblank{}.
\end{problem}

\begin{answer}
\(\dfrac{8}{5}\)
\end{answer}

\begin{problem}
双曲线 \(\dfrac{x^2}{4}-\dfrac{y^2}{5}=1\) 的右焦点到直线 \(x+2y-8=0\) 的距离为\fillinblank{}.
\end{problem}

\begin{answer}
\(\sqrt5\)
\end{answer}

\begin{problem}
记 \(\triangle ABC\) 的内角 \(A\)，\(B\)，\(C\) 的对边分别为 \(a\)，\(b\)，\(c\)，面积为 \(\sqrt3\)，\(B=60^\circ\)，\(a^2+c^2=3ac\)，则 \(b=\)\fillinblank{}.
\end{problem}

\begin{answer}
\(2\sqrt2\)
\end{answer}

\begin{problem}
以图\circled{1}为正视图，在图\circled{2}，\circled{3}，\circled{4}，\circled{5}中选两个分别作为侧视图和俯视图，组成某个三棱锥的三视图，则所选侧视图和俯视图的编号依次为\fillinblank{}（写出符合要求的一组答案即可）.

\bitmapfigure[width=17cm]{img/2021/national_paper_B_liberal/q16_fig1.png}
\end{problem}

\begin{answer}
\(\circled{2}\)\(\circled{5}\)
\end{answer}

\section{解答题}

\begin{problem}
某厂研制了一种生产高精产品的设备，为检验新设备生产产品的某项指标有无提高，用一台旧设备和一台新设备各生产了 \(10\) 件产品，得到各件产品该项指标数据如下：

\begin{center}
\begin{tabular}{|c|*{10}{c|}}
\hline
旧设备 & \(9.8\) & \(10.3\) & \(10.0\) & \(10.2\) & \(9.9\) & \(9.8\) & \(10.0\) & \(10.1\) & \(10.2\) & \(9.7\) \\
\hline
新设备 & \(10.1\) & \(10.4\) & \(10.1\) & \(10.0\) & \(10.1\) & \(10.3\) & \(10.6\) & \(10.5\) & \(10.4\) & \(10.5\) \\
\hline
\end{tabular}
\end{center}

旧设备和新设备生产产品的该项指标的样本平均数分别记为 \(\bar{x}\) 和 \(\bar{y}\)，样本方差分别记为 \(s_1^2\) 和 \(s_2^2\).

\begin{enumerate}
\item 求 \(\bar{x}\)，\(\bar{y}\)，\(s_1^2\)，\(s_2^2\)；
\item 判断新设备生产产品的该项指标的均值较旧设备是否有显著提高（如果 \(\bar{y}-\bar{x}\ge2\sqrt{\dfrac{s_1^2+s_2^2}{10}}\)，则认为新设备生产产品的该项指标的均值较旧设备有显著提高，否则不认为有显著提高）.
\end{enumerate}
\end{problem}

\solution{
\begin{enumerate}
\item 旧设备生产产品的该项指标的样本平均数和样本方差分别为
\[
\bar{x}=\frac{1}{10}\sum_{i=1}^{10} x_i=\frac{100.0}{10}=10
\]
\[
s_1^2=\frac{1}{10}\sum_{i=1}^{10}(x_i-\bar{x})^2=\frac{0.36}{10}=0.036
\]
新设备生产产品的该项指标的样本平均数和样本方差分别为
\[
\bar{y}=\frac{1}{10}\sum_{i=1}^{10} y_i=\frac{103.0}{10}=10.3
\]
\[
s_2^2=\frac{1}{10}\sum_{i=1}^{10}(y_i-\bar{y})^2=\frac{0.40}{10}=0.04
\]
\item 由（1）知 \(\bar{y}-\bar{x}=10.3-10=0.3\). 又
\[
2\sqrt{\frac{s_1^2+s_2^2}{10}}=2\sqrt{\frac{0.036+0.04}{10}}=2\sqrt{0.0076}\approx 0.174
\]
因为 \(0.3\ge 2\sqrt{0.0076}\)，即 \(\bar{y}-\bar{x}\ge2\sqrt{\dfrac{s_1^2+s_2^2}{10}}\)，所以认为新设备生产产品的该项指标的均值较旧设备有显著提高.
\end{enumerate}
}

\begin{problem}
如图，四棱锥 \(P-ABCD\) 的底面是矩形，\(PD\perp\) 底面 \(ABCD\)，\(M\) 为 \(BC\) 的中点，且 \(PB\perp AM\).

\begin{enumerate}
\item 证明：平面 \(PAM\perp\) 平面 \(PBD\)；
\item 若 \(PD=DC=1\)，求四棱锥 \(P-ABCD\) 的体积.
\end{enumerate}

\bitmapfigure[width=0.42\linewidth]{img/2021/national_paper_B_liberal/q18_fig1.png}
\end{problem}

\solution{
\begin{enumerate}
\item 因为 \(PD\perp\) 底面 \(ABCD\)，且 \(AM\subset\) 平面 \(ABCD\)，所以 \(PD\perp AM\). 又已知 \(PB\perp AM\)，且 \(PD\cap PB=P\)，所以 \(AM\perp\) 平面 \(PBD\). 因为 \(AM\subset\) 平面 \(PAM\)，所以平面 \(PAM\perp\) 平面 \(PBD\).
\item 由（1）知 \(AM\perp\) 平面 \(PBD\)，又 \(BD\subset\) 平面 \(PBD\)，所以 \(AM\perp BD\). 在矩形 \(ABCD\) 中，\(AB=DC=1\). 设 \(BC=b\)，以 \(D\) 为坐标原点，分别以 \(DC\)，\(DA\) 所在射线为 \(x\) 轴、\(y\) 轴正方向建立平面直角坐标系，则 \(D(0,0)\)，\(C(1,0)\)，\(A(0,b)\)，\(B(1,b)\). 因为 \(M\) 为 \(BC\) 的中点，所以 \(M\left(1,\dfrac{b}{2}\right)\). 则 \(\overrightarrow{AM}=\left(1,-\dfrac{b}{2}\right)\)，\(\overrightarrow{BD}=(-1,-b)\). 由 \(AM\perp BD\) 得 \(\overrightarrow{AM}\cdot\overrightarrow{BD}=-1+\dfrac{b^2}{2}=0\)，解得 \(b=\sqrt2\). 所以 \(BC=\sqrt2\)，底面面积 \(S_{ABCD}=AB\times BC=1\times\sqrt2=\sqrt2\). 故四棱锥 \(P-ABCD\) 的体积为
\[
V_{P-ABCD}=\frac13 S_{ABCD}\cdot PD=\frac13\times\sqrt2\times 1=\frac{\sqrt2}{3}
\]
\end{enumerate}
}

\begin{problem}
设 \(\{a_n\}\) 是首项为 \(1\) 的等比数列，数列 \(\{b_n\}\) 满足 \(b_n=\dfrac{na_n}{3}\). 已知 \(a_1\)，\(3a_2\)，\(9a_3\) 成等差数列.

\begin{enumerate}
\item 求 \(\{a_n\}\) 和 \(\{b_n\}\) 的通项公式；
\item 记 \(S_n\) 和 \(T_n\) 分别为 \(\{a_n\}\) 和 \(\{b_n\}\) 的前 \(n\) 项和，证明：\(T_n<\dfrac{S_n}{2}\).
\end{enumerate}
\end{problem}

\solution{
\begin{enumerate}
\item 设等比数列 \(\{a_n\}\) 的公比为 \(q\). 因为 \(a_1=1\)，且 \(a_1\)，\(3a_2\)，\(9a_3\) 成等差数列，所以
\(6a_2=a_1+9a_3\)，即 \(6q=1+9q^2\).
整理得 \((3q-1)^2=0\)，解得 \(q=\dfrac13\).
所以 \(a_n=a_1 q^{n-1}=\left(\dfrac13\right)^{n-1}\)，进而 \(b_n=\dfrac{n a_n}{3}=\dfrac{n}{3^n}\).
\item 等比数列 \(\{a_n\}\) 的前 \(n\) 项和为
\[
S_n=\frac{1-\left(\frac13\right)^n}{1-\frac13}=\frac32\left[1-\left(\frac13\right)^n\right]
\]
由 \(b_n=\dfrac{n}{3^n}\)，其前 \(n\) 项和为
\[
T_n=\frac{1}{3}+\frac{2}{3^2}+\frac{3}{3^3}+\dots+\frac{n}{3^n}
\]
两边同乘 \(\dfrac13\)，得
\[
\frac13 T_n=\frac{1}{3^2}+\frac{2}{3^3}+\dots+\frac{n-1}{3^n}+\frac{n}{3^{n+1}}
\]
两式相减得
\[
\frac23 T_n=\frac{1}{3}+\frac{1}{3^2}+\dots+\frac{1}{3^n}-\frac{n}{3^{n+1}}
=\frac{\frac13\left[1-\left(\frac13\right)^n\right]}{1-\frac13}-\frac{n}{3^{n+1}}
=\frac12\left[1-\left(\frac13\right)^n\right]-\frac{n}{3^{n+1}}
\]
两边乘 \(\dfrac32\)，得
\[
T_n=\frac34\left[1-\left(\frac13\right)^n\right]-\frac{n}{2\cdot 3^n}
\]
又 \(\dfrac{S_n}{2}=\dfrac34\left[1-\left(\dfrac13\right)^n\right]\)，且对于任意正整数 \(n\) 都有 \(\dfrac{n}{2\cdot 3^n}>0\)，所以 \(T_n<\dfrac{S_n}{2}\).
\end{enumerate}
}

\begin{problem}
已知抛物线 \(C:y^2=2px\)（\(p>0\)）的焦点 \(F\) 到准线的距离为 \(2\).

\begin{enumerate}
\item 求 \(C\) 的方程；
\item 已知 \(O\) 为坐标原点，点 \(P\) 在 \(C\) 上，点 \(Q\) 满足 \(\overrightarrow{PQ}=9\overrightarrow{QF}\)，求直线 \(OQ\) 斜率的最大值.
\end{enumerate}
\end{problem}

\solution{
\begin{enumerate}
\item 抛物线 \(C:y^2=2px\)（\(p>0\)）的焦点 \(F\left(\dfrac{p}{2},0\right)\) 到准线 \(x=-\dfrac{p}{2}\) 的距离为 \(p\). 由题设 \(p=2\)，所以 \(C\) 的方程为 \(y^2=4x\).
\item 由（1）知 \(F(1,0)\). 设点 \(P(x_0,y_0)\) 在抛物线 \(C\) 上，则 \(y_0^2=4x_0\)，即 \(x_0=\dfrac{y_0^2}{4}\). 设点 \(Q(x_Q,y_Q)\)，由 \(\overrightarrow{PQ}=9\overrightarrow{QF}\)，得
\((x_Q-x_0,y_Q-y_0)=9(1-x_Q,-y_Q)\)，
解得
\[
x_Q=\frac{x_0+9}{10}=\frac{y_0^2+36}{40},\quad y_Q=\frac{y_0}{10}
\]
显然 \(x_Q>0\). 当 \(y_0=0\) 时，\(y_Q=0\)，直线 \(OQ\) 的斜率为 \(0\). 当 \(y_0\ne 0\) 时，直线 \(OQ\) 的斜率为
\[
k=\frac{y_Q}{x_Q}=\frac{\frac{y_0}{10}}{\frac{y_0^2+36}{40}}=\frac{4y_0}{y_0^2+36}
\]
要使斜率取得最大值，必有 \(y_0>0\). 此时
\[
k=\frac{4}{y_0+\frac{36}{y_0}}\le\frac{4}{2\sqrt{y_0\cdot\frac{36}{y_0}}}=\frac{4}{12}=\frac13
\]
当且仅当 \(y_0=\dfrac{36}{y_0}\)，即 \(y_0=6\) 时等号成立. 此时 \(P(9,6)\) 在抛物线上. 所以直线 \(OQ\) 斜率的最大值为 \(\dfrac13\).
\end{enumerate}
}

\begin{problem}
已知函数 \(f(x)=x^3-x^2+ax+1\).

\begin{enumerate}
\item 讨论 \(f(x)\) 的单调性；
\item 求曲线 \(y=f(x)\) 过坐标原点的切线与曲线 \(y=f(x)\) 的公共点的坐标.
\end{enumerate}
\end{problem}

\solution{
\begin{enumerate}
\item 求导得 \(f'(x)=3x^2-2x+a\). 其判别式为 \(\Delta=4-12a=4(1-3a)\).
当 \(a\ge\dfrac13\) 时，\(\Delta\le0\)，\(f'(x)\ge0\) 恒成立，所以 \(f(x)\) 在 \(\R\) 上单调递增.
当 \(a<\dfrac13\) 时，\(\Delta>0\)，方程 \(f'(x)=0\) 有两个不相等的实数根
\[
x_1=\frac{1-\sqrt{1-3a}}{3},\quad x_2=\frac{1+\sqrt{1-3a}}{3}
\]
当 \(x\in(-\infty,x_1)\) 或 \(x\in(x_2,+\infty)\) 时，\(f'(x)>0\)；当 \(x\in(x_1,x_2)\) 时，\(f'(x)<0\). 故 \(f(x)\) 在 \(\left(-\infty,\dfrac{1-\sqrt{1-3a}}{3}\right)\) 和 \(\left(\dfrac{1+\sqrt{1-3a}}{3},+\infty\right)\) 上单调递增，在 \(\left(\dfrac{1-\sqrt{1-3a}}{3},\dfrac{1+\sqrt{1-3a}}{3}\right)\) 上单调递减.
\item 设切点坐标为 \((x_0,f(x_0))\). 则切线斜率为 \(f'(x_0)=3x_0^2-2x_0+a\)，切线方程为
\[
y-(x_0^3-x_0^2+ax_0+1)=(3x_0^2-2x_0+a)(x-x_0)
\]
因为切线过坐标原点 \((0,0)\)，所以
\[
-(x_0^3-x_0^2+ax_0+1)=(3x_0^2-2x_0+a)(-x_0)
\]
化简得 \(2x_0^3-x_0^2-1=0\)，即 \((x_0-1)(2x_0^2+x_0+1)=0\). 因为 \(2x_0^2+x_0+1=2\left(x_0+\dfrac14\right)^2+\dfrac78>0\)，所以该方程只有唯一实根 \(x_0=1\). 故切点为 \((1,1+a)\)，切线斜率为 \(k=1+a\)，切线方程为 \(y=(1+a)x\).
联立切线与曲线方程：
\[
\begin{cases}
y=(1+a)x,\\
y=x^3-x^2+ax+1,
\end{cases}
\]
得 \(x^3-x^2-x+1=0\)，即 \((x-1)^2(x+1)=0\)，解得 \(x=1\) 或 \(x=-1\).
当 \(x=1\) 时，\(y=1+a\)；当 \(x=-1\) 时，\(y=-1-a\).
所以曲线 \(y=f(x)\) 过坐标原点的切线与曲线 \(y=f(x)\) 的公共点的坐标为 \((1,1+a)\) 和 \((-1,-1-a)\).
\end{enumerate}
}

\section{选考题：解答题}

\begin{problem}
在直角坐标系 \(xOy\) 中，\(\odot C\) 的圆心为 \(C(2,1)\)，半径为 \(1\).

\begin{enumerate}
\item 写出 \(\odot C\) 的一个参数方程；
\item 过点 \(F(4,1)\) 作 \(\odot C\) 的两条切线. 以坐标原点为极点，\(x\) 轴正半轴为极轴建立坐标系，求这两条切线的极坐标方程.
\end{enumerate}
\end{problem}

\solution{
\begin{enumerate}
\item 由圆心 \(C(2,1)\)，半径为 \(1\)，可得 \(\odot C\) 的一个参数方程为
\[
\begin{cases}
x=2+\cos\theta,\\
y=1+\sin\theta,
\end{cases}
\quad\text{（\(\theta\) 为参数）}
\]
\item 点 \(F(4,1)\) 与圆心 \(C(2,1)\) 的连线平行于 \(x\) 轴，且 \(|CF|=2\). 设过点 \(F\) 的切线与直线 \(CF\) 的夹角为 \(\alpha\). 在由圆心、切点和点 \(F\) 构成的直角三角形中，\(\sin\alpha=\dfrac{1}{|CF|}=\dfrac12\)，又 \(\alpha\in\left(0,\dfrac{\pi}{2}\right)\)，所以 \(\alpha=\dfrac{\pi}{6}\). 故切线的倾斜角分别为 \(\dfrac{\pi}{6}\) 和 \(-\dfrac{\pi}{6}\).
切线方程分别为 \(y-1=\pm\dfrac{\sqrt3}{3}(x-4)\)，即
\[
x-\sqrt3y-4+\sqrt3=0,\quad x+\sqrt3y-4-\sqrt3=0
\]
将 \(x=\rho\cos\theta\)，\(y=\rho\sin\theta\) 代入第一式，得
\(\rho(\cos\theta-\sqrt3\sin\theta)=4-\sqrt3\)，即 \(2\rho\sin\left(\theta+\dfrac{5\pi}{6}\right)=4-\sqrt3\)；
代入第二式，得
\(\rho(\cos\theta+\sqrt3\sin\theta)=4+\sqrt3\)，即 \(2\rho\sin\left(\theta+\dfrac{\pi}{6}\right)=4+\sqrt3\).
所以这两条切线的极坐标方程分别为 \(2\rho\sin\left(\theta+\dfrac{5\pi}{6}\right)=4-\sqrt3\) 与 \(2\rho\sin\left(\theta+\dfrac{\pi}{6}\right)=4+\sqrt3\).
\end{enumerate}
}

\begin{problem}
已知函数 \(f(x)=|x-a|+|x+3|\).

\begin{enumerate}
\item 当 \(a=1\) 时，求不等式 \(f(x)\ge6\) 的解集；
\item 若 \(f(x)>-a\)，求 \(a\) 的取值范围.
\end{enumerate}
\end{problem}

\solution{
\begin{enumerate}
\item 当 \(a=1\) 时，\(f(x)=|x-1|+|x+3|\).
当 \(x\le -3\) 时，不等式化为 \(-(x-1)-(x+3)\ge 6\)，解得 \(x\le -4\)；
当 \(-3<x<1\) 时，不等式化为 \(-(x-1)+(x+3)\ge 6\)，即 \(4\ge 6\)，无解；
当 \(x\ge 1\) 时，不等式化为 \((x-1)+(x+3)\ge 6\)，解得 \(x\ge 2\).
综上，不等式 \(f(x)\ge 6\) 的解集为 \((-\infty,-4]\cup[2,+\infty)\).
\item 由绝对值三角不等式，对任意 \(x\in\R\)，
\[
f(x)=|x-a|+|x+3|=|a-x|+|x+3|\ge |(a-x)+(x+3)|=|a+3|
\]
当且仅当 \((a-x)(x+3)\ge 0\) 时等号成立，所以 \(f(x)\) 的最小值为 \(|a+3|\).
若 \(f(x)>-a\) 恒成立，只需 \(|a+3|>-a\).
若 \(-a<0\)，即 \(a>0\)，此时 \(|a+3|\ge 0>-a\) 恒成立；
若 \(-a\ge 0\)，即 \(a\le 0\)，此时 \(|a+3|>-a\) 等价于 \(a+3>-a\)，解得 \(a>-\dfrac32\)，结合 \(a\le 0\) 得 \(-\dfrac32<a\le 0\).
综合可知，\(a\) 的取值范围为 \(\left(-\dfrac32,+\infty\right)\).
\end{enumerate}
}
